\documentclass[a4paper]{article}

\usepackage{graphicx}
\usepackage{amsmath,amssymb}
\usepackage{minted}
\usepackage{multirow}
\usepackage{float}
\usepackage{algorithm}
\usepackage{algpseudocode}
\usepackage{todonotes}
\usepackage{forest}
\usepackage{xstring}
\usepackage{xcolor}
\usepackage[most]{tcolorbox}
\usepackage{xparse}
\usepackage{natbib}
\usepackage{adjustbox}
\usepackage[active,tightpage]{preview}

\PreviewBorder=4pt

\title{TikZ direct resize and connector interaction}
\date{2026-09-14}
\author{}

\begin{document}
\maketitle

\begin{preview}
The iPad TikZ editor now uses eight selection handles for direct node resizing,
supports in-place node-label editing, and creates directed connectors by dragging
from one node to another while the Connector tool is active. The manual Render
button was removed because the supported visual TikZ subset is rendered live on
the iPad. Source and Auto layout remain available from the options menu.

Node resize operations update position and minimum width/height in canonical
TikZ source. Connector endpoints are drawn against node boundaries and completed
connectors retain source-backed color, width, label, and arrow properties.
\end{preview}

\end{document}
